\documentclass[10pt,twoside]{extarticle}
\usepackage[paperwidth=210mm, paperheight=280mm, top=1.6cm, bottom=1.8cm, left=1.8cm, right=1.8cm]{geometry}
\usepackage{amsmath,amssymb,amsfonts}
\usepackage{graphicx}
\usepackage{tikz}
\usepackage{xcolor}
\usepackage{fancyhdr}
\usepackage{enumitem}
\usepackage{tabularx}
\usepackage[vietnamese]{babel}

% Màu sắc nhận diện Stewart Calculus
\definecolor{stewartcyan}{RGB}{0, 136, 170}
\definecolor{stewartred}{RGB}{190, 30, 45}
\definecolor{stewartdark}{RGB}{40, 40, 40}

% Biểu tượng CAS / Đồ thị
\newcommand{\casicon}{%
  \tikz[baseline=-0.6ex]{%
    \draw[stewartcyan, line width=0.6pt, rounded corners=0.8pt] (0,0) rectangle (9pt,7.5pt);%
    \draw[stewartcyan, line width=0.55pt] (1.2pt,2.2pt) sin (3.5pt,5.5pt) cos (5.5pt,1.8pt) sin (7.8pt,4.5pt);%
  }%
}

% Cấu hình Header & Footer
\pagestyle{fancy}
\fancyhf{}
\fancyhead[LE,LO]{\textbf{198} \qquad \textsc{\textbf{CHƯƠNG 3}} \quad \text{Các quy tắc tính đạo hàm}}
\fancyhead[RE,RO]{}
\renewcommand{\headrulewidth}{0pt}

\fancyfoot[C]{%
  \fontsize{6.5pt}{8pt}\selectfont\color{black!70}%
  Bản quyền thuộc về Cengage Learning. Đã đăng ký mọi quyền. Không được sao chép, quét hoặc nhân bản toàn bộ hay một phần. WCN 02-200-203\\
  Do quyền điện tử, một số nội dung của bên thứ ba có thể bị ẩn khỏi sách điện tử. Đánh giá biên tập cho thấy nội dung bị loại trừ không ảnh hưởng đến trải nghiệm học tập tổng thể. Cengage Learning có quyền xóa nội dung bổ sung vào bất kỳ lúc nào nếu các hạn chế về quyền sau này yêu cầu.
}

\setlength{\parindent}{0pt}
\setlength{\parskip}{0pt}

\begin{document}

\noindent
\begin{minipage}[t]{0.485\textwidth}
\raggedright

\textbf{23.} \ Chứng tỏ rằng $\dfrac{d}{dx}(\csc x) = -\csc x \cot x$.\par\vspace{2.5mm}

\textbf{24.} \ Chứng tỏ rằng $\dfrac{d}{dx}(\sec x) = \sec x \tan x$.\par\vspace{2.5mm}

\textbf{25.} \ Chứng tỏ rằng $\dfrac{d}{dx}(\cot x) = -\csc^2 x$.\par\vspace{2.5mm}

\textbf{26.} \ Dùng định nghĩa đạo hàm, hãy chứng minh rằng nếu $f(x) = \cos x$ thì $f'(x) = -\sin x$.\par\vspace{2.8mm}

{\color{stewartcyan}\textbf{27--30}} \ Tìm phương trình tiếp tuyến của đường cong tại điểm đã cho.\par\vspace{2mm}

\textbf{27.} \ $y = \sin x + \cos x, \quad (0, 1)$\par\vspace{2mm}

\textbf{28.} \ $y = x + \sin x, \quad (\pi, \pi)$\par\vspace{2mm}

\textbf{29.} \ $y = e^x \cos x + \sin x, \quad (0, 1)$\par\vspace{2mm}

\textbf{30.} \ $y = \dfrac{1 + \sin x}{\cos x}, \quad (\pi, -1)$\par\vspace{3.5mm}

\textbf{31.} \ \textbf{(a)} Tìm phương trình tiếp tuyến của đường cong\\
\hspace*{5.5mm}$y = 2x \sin x$ tại điểm $(\pi/2, \pi)$.\\[1.2mm]
\llap{\casicon\hspace{1.5mm}}\hspace*{5.5mm}\textbf{(b)} Minh họa phần (a) bằng cách vẽ đồ thị của đường cong và tiếp tuyến trên cùng một màn hình.\par\vspace{2.8mm}

\textbf{32.} \ \textbf{(a)} Tìm phương trình tiếp tuyến của đường cong\\
\hspace*{5.5mm}$y = 3x + 6 \cos x$ tại điểm $(\pi/3, \pi + 3)$.\\[1.2mm]
\llap{\casicon\hspace{1.5mm}}\hspace*{5.5mm}\textbf{(b)} Minh họa phần (a) bằng cách vẽ đồ thị của đường cong và tiếp tuyến trên cùng một màn hình.\par\vspace{2.8mm}

\textbf{33.} \ \textbf{(a)} Nếu $f(x) = \sec x - x$, hãy tìm $f'(x)$.\\[1.2mm]
\llap{\casicon\hspace{1.5mm}}\hspace*{5.5mm}\textbf{(b)} Kiểm tra xem câu trả lời của bạn ở phần (a) có hợp lý không bằng cách vẽ đồ thị của cả $f$ và $f'$ với $|x| < \pi/2$.\par\vspace{2.8mm}

\textbf{34.} \ \textbf{(a)} Nếu $f(x) = e^x \cos x$, hãy tìm $f'(x)$ và $f''(x)$.\\[1.2mm]
\llap{\casicon\hspace{1.5mm}}\hspace*{5.5mm}\textbf{(b)} Kiểm tra xem các câu trả lời của bạn ở phần (a) có hợp lý không bằng cách vẽ đồ thị của $f, f'$, và $f''$.\par\vspace{2.8mm}

\textbf{35.} \ Nếu $g(\theta) = \dfrac{\sin \theta}{\theta}$, hãy tìm $g'(\theta)$ và $g''(\theta)$.\par\vspace{2.8mm}

\textbf{36.} \ Nếu $f(t) = \sec t$, hãy tìm $f''(\pi/4)$.\par\vspace{2.8mm}

\textbf{37.} \ \textbf{(a)} Dùng Quy tắc thương để tính đạo hàm của hàm số
\[
f(x) = \frac{\tan x - 1}{\sec x}
\]
\hspace*{5.5mm}\textbf{(b)} Rút gọn biểu thức $f(x)$ bằng cách viết nó theo $\sin x$ và $\cos x$, rồi sau đó tìm $f'(x)$.\\[1mm]
\hspace*{5.5mm}\textbf{(c)} Chứng tỏ rằng các câu trả lời ở phần (a) và (b) là tương đương.\par\vspace{2.8mm}

\textbf{38.} \ Giả sử $f(\pi/3) = 4$ và $f'(\pi/3) = -2$, và đặt\\
$g(x) = f(x)\sin x$ và $h(x) = (\cos x)/f(x)$. Tìm\\[1mm]
\hspace*{5.5mm}\textbf{(a)} $g'(\pi/3)$ \hfill \textbf{(b)} $h'(\pi/3)$\hspace*{15mm}\null\par\vspace{3mm}

{\color{stewartcyan}\textbf{39--40}} \ Với những giá trị nào của $x$ thì đồ thị của $f$ có tiếp tuyến nằm ngang?\par\vspace{2mm}

\textbf{39.} \ $f(x) = x + 2\sin x$ \hfill \textbf{40.} \ $f(x) = e^x \cos x$

\end{minipage}%
\hfill
\begin{minipage}[t]{0.485\textwidth}
\raggedright

\textbf{41.} \ Một vật gắn vào lò xo dao động theo phương ngang trên một mặt phẳng nằm ngang nhẵn (xem hình). Phương trình chuyển động của nó là $x(t) = 8 \sin t$, trong đó $t$ tính bằng giây và $x$ tính bằng xentimét.
\begin{enumerate}[label=\textbf{(\alph*)}, leftmargin=6mm, itemsep=0.8mm, topsep=1mm]
  \item Tìm vận tốc và gia tốc tại thời điểm $t$.
  \item Tìm vị trí, vận tốc và gia tốc của vật tại thời điểm $t = 2\pi/3$. Tại thời điểm đó vật đang chuyển động theo hướng nào?
\end{enumerate}

\begin{center}
  \vspace{-0.5mm}
  \includegraphics[width=0.74\linewidth]{images/p0233_fig1.png}
  \vspace{-1.5mm}
\end{center}

\llap{\casicon\hspace{1.5mm}}\textbf{42.} \ Một dây đàn hồi được treo vào một chiếc móc và một vật nặng được treo vào đầu dưới của dây. Khi vật bị kéo xuống phía dưới rồi thả ra, nó dao động theo phương thẳng đứng. Phương trình chuyển động là $s = 2\cos t + 3\sin t$, $t \geqslant 0$, trong đó $s$ được đo bằng xentimét và $t$ bằng giây. (Chọn chiều dương hướng xuống dưới.)
\begin{enumerate}[label=\textbf{(\alph*)}, leftmargin=6mm, itemsep=0.6mm, topsep=1mm]
  \item Tìm vận tốc và gia tốc tại thời điểm $t$.
  \item Vẽ đồ thị các hàm vận tốc và gia tốc.
  \item Khi nào thì vật đi qua vị trí cân bằng lần đầu tiên?
  \item Vật di chuyển cách vị trí cân bằng xa nhất là bao nhiêu?
  \item Khi nào thì tốc độ đạt giá trị lớn nhất?
\end{enumerate}\par\vspace{2mm}

\textbf{43.} \ Một chiếc thang dài $10\text{ ft}$ tựa vào một bức tường thẳng đứng. Gọi $\theta$ là góc giữa đỉnh thang và tường, và $x$ là khoảng cách từ chân thang đến tường. Nếu chân thang trượt ra xa khỏi tường, hỏi $x$ thay đổi nhanh như thế nào đối với $\theta$ khi $\theta = \pi/3$?\par\vspace{2.5mm}

\textbf{44.} \ Một vật có trọng lượng $W$ được kéo dọc theo một mặt phẳng nằm ngang nhờ một lực tác dụng dọc theo sợi dây buộc vào vật. Nếu sợi dây tạo với mặt phẳng một góc $\theta$, thì độ lớn của lực kéo là
\[
F = \frac{\mu W}{\mu \sin \theta + \cos \theta}
\]
trong đó $\mu$ là một hằng số được gọi là \emph{hệ số ma sát}.
\begin{enumerate}[label=\textbf{(\alph*)}, leftmargin=6mm, itemsep=0.8mm, topsep=1mm]
  \item Tìm tốc độ thay đổi của $F$ theo $\theta$.
  \item Khi nào thì tốc độ thay đổi này bằng $0$?
  \item[\llap{\casicon\hspace{1.5mm}}\textbf{(c)}] Nếu $W = 50\text{ lb}$ và $\mu = 0{,}6$, hãy vẽ đồ thị của $F$ như một hàm số của $\theta$ và sử dụng đồ thị đó để định vị giá trị của $\theta$ mà tại đó $dF/d\theta = 0$. Giá trị này có nhất quán với câu trả lời của bạn ở phần (b) không?
\end{enumerate}\par\vspace{3mm}

{\color{stewartcyan}\textbf{45--60}} \ Tìm giới hạn.\par\vspace{2.5mm}

\noindent
\begin{tabularx}{\linewidth}{@{}X X@{}}
\textbf{45.} $\displaystyle \lim_{x \to 0} \frac{\sin 5x}{3x}$ & \textbf{46.} $\displaystyle \lim_{x \to 0} \frac{\sin x}{\sin \pi x}$ \\[3.8ex]
\textbf{47.} $\displaystyle \lim_{t \to 0} \frac{\sin 3t}{\sin t}$ & \textbf{48.} $\displaystyle \lim_{x \to 0} \frac{\sin^2 3x}{x}$
\end{tabularx}

\end{minipage}

\end{document}